% =============================================================================
% Theory
% =============================================================================
\section{Theory}\label{sec:theory}

\subsection{Equivalent-sphere forward model}

We use the flexural-frequency model introduced in the companion
paper~\cite{BrowntoneCompanion}, drawing on classical shell vibration
theory~\cite{Soedel2004,Leissa1973}.  The shell has semi-major axis $a$,
semi-minor axis $c$, thickness $h$, Young's modulus $E$, Poisson's ratio $\nu$,
wall density $\rho_w$, and internal fluid density $\rho_f$.  The fluid bulk modulus
is denoted by $K_f$, the static internal pressure by $P_\mathrm{iap}$, and the
loss tangent by $\eta$.  Unless stated otherwise, the canonical parameter set is
\[
a=\SI{0.18}{m},\quad
c=\SI{0.12}{m},\quad
h=\SI{0.01}{m},\quad
E=\SI{0.1}{MPa},\quad
\nu=0.45,
\]
\[
\rho_w=\SI{1100}{kg/m^3},\quad
\rho_f=\SI{1020}{kg/m^3},\quad
K_f=\SI{2.2}{GPa},\quad
P_\mathrm{iap}=\SI{1000}{Pa},\quad
\eta=0.25.
\]

The model collapses the geometry to the equivalent radius
\begin{equation}\label{eq:req}
  R_\mathrm{eq} = (a^2 c)^{1/3}.
\end{equation}
For flexural modes $n \geq 2$, the angular frequency $\omega_n=2\pi f_n$ is
written as
\begin{equation}\label{eq:flexural_forward}
  \omega_n^2
  =
  \frac{K_\mathrm{bend}^{(n)} + K_\mathrm{memb}^{(n)} + K_\mathrm{pre}^{(n)}}
       {m_\mathrm{eff}^{(n)}},
\end{equation}
where the effective stiffness contains bending, membrane, and prestress terms,
and the effective mass contains wall mass and fluid added mass.

With flexural rigidity
\begin{equation}
  D = \frac{E h^3}{12(1-\nu^2)},
\end{equation}
the individual contributions are
\begin{equation}
  K_\mathrm{bend}^{(n)}
  =
  \frac{n(n-1)(n+2)^2}{R_\mathrm{eq}^4}\,D,
\end{equation}
\begin{equation}
  K_\mathrm{memb}^{(n)}
  =
  \frac{E h}{R_\mathrm{eq}^2}\,\lambda_n,
  \qquad
  \lambda_n
  =
  \frac{n^2+n-2+2\nu}{n^2+n+1-\nu},
\end{equation}
\begin{equation}
  K_\mathrm{pre}^{(n)}
  =
  \frac{P_\mathrm{iap}}{R_\mathrm{eq}}(n-1)(n+2),
\end{equation}
and
\begin{equation}
  m_\mathrm{eff}^{(n)}
  =
  \rho_w h + \frac{\rho_f R_\mathrm{eq}}{n}.
\end{equation}

Equation~\eqref{eq:flexural_forward} defines the equivalent-sphere reduction
used as a degenerate comparator throughout the paper.  It depends on $(a,c)$
only through $R_\mathrm{eq}$ and is therefore algebraically rank deficient
(Proposition~\ref{prop:chain_rule}).  The oblate Ritz model introduced in
Section~\ref{sec:perturbation} is the forward model used for all main results.

\subsection{Inverse problem}

Let
\begin{equation}
  \boldsymbol{\theta} = (a,c,E)^\top
\end{equation}
denote the parameter vector to be estimated, and let
\begin{equation}
  \mathbf{f}(\boldsymbol{\theta})
  =
  \bigl(f_2(\boldsymbol{\theta}),
        f_3(\boldsymbol{\theta}),
        \ldots,
        f_N(\boldsymbol{\theta})\bigr)^\top
\end{equation}
be the vector of measured flexural frequencies.  The inverse problem is to
recover $\boldsymbol{\theta}$ from the observation
$\mathbf{f}^\mathrm{obs}$.

Local identifiability is controlled by the Jacobian
\begin{equation}\label{eq:jacobian_raw}
  J_{ij}
  =
  \frac{\partial f_i}{\partial \theta_j},
\end{equation}
where $i$ indexes the measured mode and $j$ indexes the unknown parameter.
Because the raw Jacobian mixes units of \si{Hz/m} and \si{Hz/Pa}, its condition
number is not a meaningful identifiability metric by itself.  We therefore use
the dimensionless scaled Jacobian
\begin{equation}\label{eq:jacobian_scaled}
  J_s[i,j]
  =
  \left(\frac{\partial f_i}{\partial \theta_j}\right)
  \left(\frac{\theta_j}{f_i}\right),
\end{equation}
which measures fractional frequency change per fractional parameter change.

The singular values $\sigma_1 \ge \sigma_2 \ge \sigma_3$ of $\mathbf{J}_s$
quantify the observable directions in parameter space, and the corresponding
condition number
\begin{equation}\label{eq:kappa}
  \kappa(\mathbf{J}_s)
  =
  \frac{\sigma_{\max}}{\sigma_{\min}}
\end{equation}
measures the worst-case amplification of small frequency errors into parameter
errors.  A large $\kappa$ indicates a nearly unobservable parameter
combination; a divergent $\kappa$ indicates rank deficiency.

\subsection{Fisher information and Cram\'er--Rao bounds}

Assume independent fractional frequency errors with standard deviation
$\epsilon_\mathrm{noise}$.  In dimensionless form, the Fisher information
matrix becomes
\begin{equation}\label{eq:fisher}
  \mathbf{F}
  =
  \frac{1}{\epsilon_\mathrm{noise}^2}\,
  \mathbf{J}_s^\top \mathbf{J}_s.
\end{equation}
Equivalently, one may interpret \eqref{eq:fisher} as the Fisher matrix of the
whitened residuals.  The Cram\'er--Rao lower bound then states that the
covariance of any unbiased local estimator satisfies
\begin{equation}\label{eq:crlb}
  \mathrm{Cov}(\hat{\boldsymbol{\theta}})
  \succeq
  \mathbf{F}^{-1},
\end{equation}
with the diagonal entries of $\mathbf{F}^{-1}$ giving the minimum achievable
variances under the assumed noise model.

This framework does not prove global uniqueness.  It answers a more practical
question: at a chosen operating point, do the measured resonances contain enough
independent information to recover $(a,c,E)$ with tolerable sensitivity to
noise?  That is a question of local practical identifiability, and it is the
question addressed in the next sections.

\subsection{Analytical rank deficiency of the equivalent-sphere Jacobian}

Before computing any Jacobian numerically, it is useful to understand the
spherical rank deficiency from first principles.  The following result makes the
mechanism explicit.

\begin{proposition}[Chain-rule rank deficiency]\label{prop:chain_rule}
For any forward model in which all frequencies $f_n$ depend on the geometric
pair $(a,c)$ only through the equivalent radius
$R_\mathrm{eq} = (a^2 c)^{1/3}$, the geometric sub-Jacobian
$[\partial \mathbf{f}/\partial a,\;\partial \mathbf{f}/\partial c]$ has rank
one.  Consequently, the full Jacobian
$\mathbf{J} = [\partial \mathbf{f}/\partial a,\;
\partial \mathbf{f}/\partial c,\;
\partial \mathbf{f}/\partial E]$
has rank at most two (one fewer than the three parameter directions), and
the scaled Jacobian $\mathbf{J}_s$ inherits the same rank deficiency.
\end{proposition}

\begin{pf}
By the chain rule, the partial derivatives of any frequency $f_n$ with respect
to $a$ and $c$ factor through $R_\mathrm{eq}$:
\begin{equation}\label{eq:dreq_da}
  \frac{\partial R_\mathrm{eq}}{\partial a}
  = \frac{2}{3}\,\frac{R_\mathrm{eq}}{a},
  \qquad
  \frac{\partial R_\mathrm{eq}}{\partial c}
  = \frac{1}{3}\,\frac{R_\mathrm{eq}}{c}.
\end{equation}
Hence
\begin{equation}\label{eq:dfn_chain}
  \frac{\partial f_n}{\partial a}
  = \frac{\mathrm{d} f_n}{\mathrm{d} R_\mathrm{eq}}
    \cdot \frac{2\,R_\mathrm{eq}}{3a},
  \qquad
  \frac{\partial f_n}{\partial c}
  = \frac{\mathrm{d} f_n}{\mathrm{d} R_\mathrm{eq}}
    \cdot \frac{R_\mathrm{eq}}{3c}.
\end{equation}
The ratio of the two partial derivatives is therefore
\begin{equation}\label{eq:ratio_chain}
  \frac{\partial f_n / \partial a}
       {\partial f_n / \partial c}
  = \frac{2c}{a}
  \qquad \text{for all } n,
\end{equation}
independently of the functional form of $f_n(R_\mathrm{eq})$.  The
corresponding scaled-Jacobian columns satisfy
\begin{equation}
  J_s[i,a] \;=\; \frac{\partial f_i}{\partial a}\,\frac{a}{f_i}
  \;=\; 2\,\frac{\partial f_i}{\partial c}\,\frac{c}{f_i}
  \;=\; 2\,J_s[i,c].
\end{equation}
Two columns of $\mathbf{J}_s$ are therefore proportional (with constant of
proportionality $2$, independent of mode index $i$), so the geometric
sub-Jacobian has rank one and the full $\mathbf{J}_s$ has rank at most two.
\end{pf}

\begin{remark}[Value reduction versus directional-derivative structure]
The rank deficiency at $\varepsilon = 0$ has two distinct aspects that should
not be conflated.  First, the \emph{values} of the frequencies become
degenerate: modes that would be distinguishable on an oblate shell cluster
together as the geometry approaches a sphere, reducing the effective
information content of the spectrum.  Second, the \emph{directional-derivative
structure} of the Jacobian degenerates: the columns for $a$ and $c$ become
exactly proportional (by~\eqref{eq:ratio_chain}), rendering a parameter
combination unobservable regardless of spectral precision.  Both effects
worsen identifiability, but the second is the more fundamental obstruction
--- it is algebraic rather than numerical and persists even with perfect
frequency measurements.
\end{remark}

\begin{remark}
The ratio~\eqref{eq:ratio_chain} is purely geometric: it reflects the exponents
$(2,1)$ in $R_\mathrm{eq} = (a^2 c)^{1/3}$ and holds regardless of whether the
forward model includes bending, membrane, prestress, or added-mass terms.  Any
model built on $R_\mathrm{eq}$ alone inherits this rank deficiency.
\end{remark}

\subsection{Eccentricity perturbation and identifiability restoration}\label{sec:perturbation}

Proposition~\ref{prop:chain_rule} identifies the root cause of
non-identifiability at the spherical limit.  The natural follow-up is:
\emph{how} does oblate geometry restore full rank?  The answer lies in
mode-dependent curvature sampling.

For an oblate spheroid with semi-axes $(a,c)$, define the eccentricity
\begin{equation}\label{eq:eccentricity}
  \varepsilon = \sqrt{1 - c^2/a^2}.
\end{equation}
At $\varepsilon = 0$ the shell is a sphere; as $\varepsilon \to 1$ it becomes a
thin disc.

\paragraph{Oblate Ritz forward model.}
The oblate Ritz model retains the same physical stiffness contributions as
the equivalent-sphere model but integrates them over the true spheroidal
geometry.  The trial displacement is expanded in Legendre modes
$P_n(\cos\theta)$ for $n \in \{2,3,\ldots,N\}$, with each mode treated as
an independent generalised coordinate $q_n$.  The resulting Ritz eigenvalue
problem for mode $n$ is
\begin{equation}\label{eq:ritz_eigenvalue}
  \omega_n^2
  =
  \frac{\tilde K_\mathrm{bend}^{(n)}
      + \tilde K_\mathrm{memb}^{(n)}
      + \tilde K_\mathrm{pre}^{(n)}}
       {\tilde m_\mathrm{eff}^{(n)}},
\end{equation}
where the tildes distinguish Ritz stiffness and mass from their
equivalent-sphere counterparts.  Specifically, the bending stiffness is
\begin{equation}\label{eq:ritz_bend}
  \tilde K_\mathrm{bend}^{(n)}
  =
  D \int_0^1
    \frac{[n(n+1)]^2\,P_n^2(\eta)}{G(\eta)^2}
  \,\mathrm{d}\eta,
\end{equation}
the membrane stiffness is
\begin{equation}\label{eq:ritz_memb}
  \tilde K_\mathrm{memb}^{(n)}
  =
  Eh \int_0^1
    \frac{\lambda_n\,P_n^2(\eta)}{G(\eta)}
  \,\mathrm{d}\eta,
\end{equation}
and the effective mass is
\begin{equation}\label{eq:ritz_mass}
  \tilde m_\mathrm{eff}^{(n)}
  =
  \int_0^1
    \left(\rho_w h + \frac{\rho_f\,\sqrt{G(\eta)}}{n}\right)
    P_n^2(\eta)
  \,\mathrm{d}\eta,
\end{equation}
with the geometric factor $G(\eta) = a^2\eta^2 + c^2(1-\eta^2)$ encoding
meridional curvature variation.  The prestress contribution from
intra-abdominal pressure is
\begin{equation}\label{eq:ritz_pre}
  \tilde K_\mathrm{pre}^{(n)}
  =
  P_\mathrm{iap} \int_0^1
    N_1(\eta)\,\frac{(1-\eta^2)\,[P_n'(\eta)]^2}{G(\eta)}
  \,\sqrt{G(\eta)}\,\mathrm{d}\eta,
\end{equation}
where $N_1(\eta)$ is the meridional membrane stress resultant under
uniform internal pressure.
In the spherical limit $a = c$, the integrands become
$\eta$-independent and the schematic integrals above recover the functional
structure of the equivalent-sphere expressions, though the truncated
$2\times2$ Ritz eigenproblem incurs a systematic offset of approximately
$8\%$ in eigenfrequencies relative to the analytical sphere formula.  This
offset does not affect the identifiability analysis, which depends on
sensitivity ratios (Jacobian column directions) rather than absolute
frequency values.  Near the spherical limit ($c/a > 0.98$),
the oblate fluid harmonics are blended to sphere-approximated values using a
$C^2$-continuous smootherstep function over $c/a \in [0.98,\,0.999]$, ensuring
numerical stability without discontinuous derivatives.
For $a \neq c$, each mode
``sees'' a different effective radius through the $G(\eta)$ weighting,
which is the mechanism by which eccentricity breaks the rank degeneracy.

The expressions~\eqref{eq:ritz_eigenvalue}--\eqref{eq:ritz_pre} are
schematic reductions of the Ritz quotient, presented for physical
interpretation of the individual stiffness and mass contributions.  The
implemented Ritz solver constructs, for each harmonic order~$n$, a $2\times2$
generalised eigenproblem $\mathbf{K}^{(n)}\mathbf{q}
= \omega^2\mathbf{M}^{(n)}\mathbf{q}$ coupling normal ($\beta$) and
tangential ($\alpha$) displacement coordinates.  The stiffness matrix
$\mathbf{K}^{(n)}$ includes bending, membrane, prestress, and
cross-coupling contributions; the mass matrix $\mathbf{M}^{(n)}$ includes
shell and added fluid mass terms.  The reported eigenfrequency is the
lowest eigenvalue branch of this coupled problem.  Full implementation
details and source code are provided in the companion repository.

\begin{remark}[Mode-dependent curvature sampling]\label{prop:eccentricity}
The geometric factor $G(\eta)$ enters the Ritz stiffness and mass integrals
nonlinearly.  Because the oblate spheroid has reflection symmetry about the
equatorial plane, odd powers of $\varepsilon$ vanish in any expansion of the
eigenvalues or their sensitivities, so the leading correction to
$\sigma_{\min}$ is $O(\varepsilon^2)$.  This is the standard parity argument
of analytic perturbation theory~\cite{Kato1966,StewartSun1990}: for a
one-parameter family $A(\varepsilon)$ with $A(\varepsilon) = A(-\varepsilon)$,
the eigenvalue branches are even functions of $\varepsilon$.

The physical mechanism can be described as follows.
\begin{enumerate}
  \item The stiffness integral for mode $n$,
  \begin{equation}\label{eq:Ink}
    I_n^{(k)} = \int_0^1 P_n^2(\eta)\, G^{-k}(\eta)\,\mathrm{d}\eta,
  \end{equation}
  acquires $n$-dependent corrections at order $\varepsilon^2$, because
  $P_n^2(\eta)$ samples the interval $[0,1]$ differently for each mode number.

  \item Consequently, the partial derivatives $\partial f_n / \partial a$ and
  $\partial f_n / \partial c$ acquire $n$-dependent corrections:
  \begin{equation}\label{eq:ratio_corrected}
    \frac{\partial f_n / \partial a}
         {\partial f_n / \partial c}
    = \frac{2c}{a} + \varepsilon^2\,\delta_n + O(\varepsilon^4),
  \end{equation}
  where $\delta_n$ depends on the mode number through curvature-weighted
  integrals of $P_n^2$.

  \item Since the ratio is no longer constant across modes, the columns for $a$
  and $c$ in the Jacobian are no longer proportional, and
  $\mathbf{J}_s$ acquires full rank for $\varepsilon > 0$.
\end{enumerate}
The $O(\varepsilon^2)$ scaling and the numerical values of the expansion
coefficients are established by polynomial regression in
Section~\ref{sec:results}; we do not claim a closed-form analytical prefactor.
\end{remark}

The mechanism is \emph{mode-dependent curvature sampling}: different
Legendre modes $P_n(\eta)$ weight the meridional curvature gradient
$\partial G / \partial \eta$ differently.  Near the poles ($\eta \to 1$) the
oblate shell is flatter; near the equator ($\eta \to 0$) it is more curved.
Low-order modes ($n=2$) sample both regions broadly, while higher-order modes
($n \geq 4$) concentrate energy near the equator.  This differential sampling
breaks the proportionality of the $a$- and $c$-columns and lifts the rank
deficiency.

Over the intermediate eccentricity range $\varepsilon \in [0.15,\,0.85]$,
the minimum singular value is well described by the polynomial fit
\begin{equation}\label{eq:sigma_min_expansion}
  \sigma_{\min}(\varepsilon)
  \approx
  \sigma_0 + \lambda_1\,\varepsilon^{2} + \lambda_2\,\varepsilon^{4},
\end{equation}
where $\sigma_0$, $\lambda_1$, $\lambda_2$ are regression coefficients
(not asymptotic limits).  The expansion is empirical: $\sigma_0$ is the
regression intercept extrapolated to $\varepsilon = 0$, not the directly
computed value at the sphere.  Direct evaluation at $\varepsilon = 0.01$
(five modes) gives $\kappa_\mathrm{floor} \approx 269$, supporting a finite
near-sphere $\kappa$ for the implemented model rather than divergence.
The crossover eccentricity
$\varepsilon_c = \sqrt{\sigma_0/\lambda_1}$ separates a regime dominated
by the regression intercept ($\varepsilon < \varepsilon_c$) from one
dominated by the eccentricity-driven terms ($\varepsilon > \varepsilon_c$).
Section~\ref{sec:results} determines
$\sigma_0$, $\lambda_1$, $\lambda_2$, and $\varepsilon_c$ numerically and
confirms that for the oblate Ritz model, $\varepsilon_c > 1$: the intercept
dominates across the entire physical eccentricity range.
